\section{The exact receiving map into the original mixed factors}
\label{sec:OPR}

The current receiver MRI1--5 is retained as a frozen predecessor.
This section proves its next observation input from OPG1--26.
Fix either original source order $s\in\{1,k\}$, put $D=q+j$,
and retain all the fixed coefficient maps of OPG and GMB.
In particular $e_D=[u_D]$, $\lambda_D=\Lambda e_D$, and
$\mathcal M_D=\sum_{d<D}\lambda_d\lambda_d^*$ on the original
boundary image. At the preceding polynomial degree $D-1$, GMB4
and OPG15 give
\[
\begin{gathered}
 G_{D-1}^{(s)}=C_s^{-*}K_j^{-1}C_s^{-1},\quad
 (G_{D-1}^{(s)})^{-1}=C_sK_jC_s^*,\\
 Q_{D-1}^{(s)}=\mathcal M_D^{-1},\qquad
 S_{D-1}^{(s)}=C_sK_jC_s^*\Lambda^*\mathcal M_D^{-1}:B\to E.
\end{gathered}\tag{OPR1}
\]
Indeed $\Lambda(G_{D-1}^{(s)})^{-1}\Lambda^*=
R_sK_jR_s^*=\mathcal M_D$. Direct multiplication therefore
proves $\Lambda S_{D-1}^{(s)}=I_B$ and
$I^*G_{D-1}^{(s)}S_{D-1}^{(s)}=0$. Thus this is precisely
the original minimum section for this degree and this source.
All statements use their unique empty meaning if $B=0$.

MRI2 uses the kernel vector
$r_j^{\rm ker}=v_j-K_jR_s^*\mathcal M_D^{-1}\lambda_D$
in the $C_s$ coordinates. The fixed-section kernel vector of
OPG21 is $\zeta_D=\kappa e_D$, with $I\kappa=I_E-S_*\Lambda$.
The exact morphism between these two specified vectors is
\[
\boxed{\begin{aligned}
 C_s r_j^{\rm ker}
   &=e_D-S_{D-1}^{(s)}\lambda_D\\
   &=I\zeta_D+(S_*-S_{D-1}^{(s)})\lambda_D
    =I\vartheta_{D,j}^{(s)},\\
 \vartheta_{D,j}^{(s)}
   &=\zeta_D-\kappa S_{D-1}^{(s)}\lambda_D\in K.
\end{aligned}}\tag{OPR2}
\]
For the first equality substitute $v_j=C_s^{-1}e_D$ and
$R_s^*=C_s^*\Lambda^*$ into MRI2. For the second substitute
$e_D=I\zeta_D+S_*\lambda_D$. The difference of the two
sections lands in $K$ because their observations are both $I_B$.
Applying $\kappa$ and using $\kappa S_*=0$ proves the last
line and hence the third equality. This proves the complete
section correction in the original coefficient space.

Retain the actual kernel Gram
$H_{D-1}^{(s),K}=I^*G_{D-1}^{(s)}I$. Congruencing OPR1 by
$C_s$ and then using OPR2 proves
\[
 t_j-\xi_j
 = (r_j^{\rm ker})^*K_j^{-1}r_j^{\rm ker}
 = (\vartheta_{D,j}^{(s)})^*
                    H_{D-1}^{(s),K}\vartheta_{D,j}^{(s)}.
 \tag{OPR3}
\]
The first equality is the proved MRI2/GMB10 orthogonal
decomposition. Thus the recurrence OPG21 supplies both the
actual boundary column and, through the complete section
correction OPR2, the exact kernel norm needed at each step.
The original ordered factors are consequently
\[
\begin{gathered}
 \xi_j=\lambda_D^*\mathcal M_D^{-1}\lambda_D,\qquad
 1+\frac{\nu_{D-1}}{a_{D-1}+\beta_{D-1}}=1+\xi_j,\\
 1+\frac{\beta_{D-1}}{a_{D-1}}
 =\frac{1+t_j}{1+\xi_j}
 =1+\frac{(\vartheta_{D,j}^{(s)})^*
                  H_{D-1}^{(s),K}\vartheta_{D,j}^{(s)}}{1+\xi_j}.
\end{gathered}\tag{OPR4}
\]
The first line uses OPG17 and GMB13. The second line follows
by inserting OPR3 into that same GMB13 identity. This preserves
the denominator $1+\xi_j$ and the original order of elimination.

The finite endpoint operation remains
$\mathcal W_q(f)=f_0+2\sum_{j=1}^{q-1}f_j+f_q$.
Applying it to the logarithms in OPR4 reproduces MRI4 on its
unchanged source. The exact arithmetic receiving equality is
therefore still
\[
\begin{aligned}
 \mathcal B_k^{\rm ar}
 &=F_K^{(k)}+F_B^{(k)}+\mathcal H_k+\delta_k^\sigma\\
 &=F_K^{(1)}+F_B^{(1)}+\delta_k^\sigma.
\end{aligned}\tag{OPR5}
\]
Its proof is substitution into the existing MRI5/BRI2 source
identity: $F_K^{(k)}+F_B^{(k)}=\mathcal B_k^0$,
$F_K^{(1)}+F_B^{(1)}=\mathcal B_k^\sigma$,
$\mathcal H_k=\mathcal B_k^\sigma-\mathcal B_k^0$, and
$\delta_k^\sigma=\mathcal B_k^{\rm ar}-\mathcal B_k^\sigma$.
No term is assigned a new sign or asymptotic coefficient.
The combined coefficient received by BRI6 and BRI9 remains
attached to this same sum. OPR1--4 provide its individual
original factors from the actual OPG observation vectors,
with all original source transitions present in OPR5.
